\documentclass[12pt,a4paper]{article}
\usepackage[UTF8]{ctex}
\usepackage{amsmath,amssymb}
\usepackage{geometry}
\usepackage{booktabs}
\usepackage{array}
\usepackage{enumitem}
\usepackage{hyperref}

\geometry{left=2.5cm,right=2.5cm,top=2.5cm,bottom=2.5cm}

\title{束语义（Sheaf Semantics）}
\author{}
\date{}

\begin{document}

\maketitle

\section{概述与核心构想}

束语义（Sheaf Semantics），又称 Kripke-Joyal 语义，是范畴论、数理逻辑与拓扑学交叉领域中的核心工具。

在经典逻辑与模型论（如 Tarski 语义）中，一个命题的真值只能是“真”或“假”（即二值集合 $\{0, 1\}$）。但在拓扑与几何语境下，数学对象的性质往往随“局部区域/参数”改变。

束语义的思想是将真值空间从二值集合扩展为某个空间（或 Locale）$X$ 的开集格 $\mathcal{O}(X)$。一个命题 $\varphi$ 在开集 $U \in \mathcal{O}(X)$ 上成立，记作：

\[
U \models \varphi
\]

读作 “开集 $U$ 迫近（Force）或验证了命题 $\varphi$”。

束语义的核心价值在于：它允许我们将在空间 $X$ 上连续变化的对象（即束 Sheaves）看作普通集合，并在束的范畴内部使用直觉主义一阶逻辑（Intuitionistic First-Order Logic）进行推演。

\section{预备概念与符号定义}

为了给出无歧义且逻辑清晰的解释，下文严谨定义了束语义所依赖的所有底层概念。

\subsection{2.1 偏序集与格 (Poset and Lattice)}

\begin{itemize}
  \item 偏序集 (Poset)：一个二元组 $(P, \le)$，其中 $P$ 为集合，$\le$ 是自反、反对称和传递的自洽二元关系。
  \item 交 (Meet $\wedge$) 与并 (Join $\bigvee$)：若 $P$ 中任意有限子集都有最大下界，称为有限交（Meet $\wedge$）；若任意子集（包含无限集）都有最小上界，称为任意并（Join $\bigvee$）。
  \item 最大元与最小元：空交记作最大元 $\top$（Top），空并记作最小元 $\perp$（Bottom）。
\end{itemize}

\subsection{2.2 框架与局部拓扑 (Frame \& Locale)}

\begin{itemize}
  \item 框架 (Frame)：一个偏序集 $L$，若它满足以下条件：
  \begin{enumerate}
    \item 拥有任意并 $\bigvee$ 和有限交 $\wedge$。
    \item 满足无限分配律（Infinitary Distributive Law）：对于任意 $U \in L$ 以及任意子集族 $\{V_i\}_{i \in I} \subseteq L$，均有：
    \[
    U \wedge \left( \bigvee_{i \in I} V_i \right) = \bigvee_{i \in I} (U \wedge V_i)
    \]
  \end{enumerate}
  \item 框架态射 (Frame Homomorphism)：单调映射 $\alpha: L \to L'$，且保持任意并和有限交。
  \item 局部拓扑 (Locale)：Locale 是 Frame 的范畴对偶概念。一个 Locale $X$ 由其“开集框架” $\mathcal{O}(X)$ 决定。Locale 之间的态射 $f: X \to X'$ 对应反向的框架态射 $f^*: \mathcal{O}(X') \to \mathcal{O}(X)$。
  \item 开集的包含关系：在 $\mathcal{O}(X)$ 中，记 $V \le U$ 表示开集 $V$ 包含于开集 $U$。
\end{itemize}

\subsection{2.3 预束与束 (Presheaf and Sheaf)}

设 $X$ 为一个 Locale（或拓扑空间）。

\begin{itemize}
  \item 预束 (Presheaf)：Locale $X$ 上的预束 $F$ 是一个反向函子 $F: \mathcal{O}(X)^{\text{op}} \to \mathbf{Set}$，具体包含：
  \begin{enumerate}
    \item 对每个开集 $U \in \mathcal{O}(X)$，赋予一个集合 $F(U)$（其元素称为 $F$ 在 $U$ 上的截面 Section）。
    \item 对任意 $V \le U$，赋予一个限制映射（Restriction Map） $r_V^U: F(U) \to F(V)$，通常记作 $s\vert{}_V := r_V^U(s)$（对于 $s \in F(U)$）。
    \item 满足自洽性：$s\vert{}_U = s$ 且对 $W \le V \le U$，有 $(s\vert{}_V)\vert{}_W = s\vert{}_W$。
  \end{enumerate}
  \item 相容截面族 (Compatible Family)：若开集 $U$ 被开集族覆盖，即 $U = \bigvee_{i \in I} U_i$，而截面族 $(s_i)_{i \in I}$ 满足 $s_i \in F(U_i)$，且对任意 $i, j \in I$，在交集 $U_i \wedge U_j$ 上约束一致：
  \[
  s_i\vert{}_{U_i \wedge U_j} = s_j\vert{}_{U_i \wedge U_j}
  \]
  则称此族为相容族。
  \item 束 (Sheaf)：若预束 $F$ 对任意覆盖 $U = \bigvee_{i \in I} U_i$ 及任意相容截面族 $(s_i)_{i \in I}$，都存在唯一截面 $s \in F(U)$ 使得对所有 $i \in I$ 均有 $s\vert{}_{U_i} = s_i$，则称 $F$ 为 Locale $X$ 上的束。
  \item 束态射 (Morphism of Sheaves)：束 $F$ 到 $G$ 的自然变换 $f: F \to G$，即由一组映射族 $f_U: F(U) \to G(U)$ 组成，且与限制映射可交换。
\end{itemize}

\section{一阶语言与语法结构 (Language \& Syntax)}

为了定义束 $F, G, \dots$ 上的逻辑语义，我们需要建立基于 Locale $X$ 的一阶类型化语言 (First-Order Typed Language)：

\begin{enumerate}
  \item 类型/排序 (Sorts)：每个束 $F$ 作为一个 Sort（类型）。
  \item 项与变量 (Terms and Variables)：在开集 $U \in \mathcal{O}(X)$ 上，Sort $F$ 的项/常量是 $F(U)$ 中的截面 $s \in F(U)$。
  \item 函数符号 (Function Symbols)：束态射 $f: F \to G$ 视为作用于截面的函数。
  \item 公式 (Formulas)：由等式 $s = t$、逻辑常值 $\top, \perp$、逻辑连结词 $\wedge, \vee, \Rightarrow$ 以及量词 $\forall, \exists$ 递归构建。
\end{enumerate}

\section{束语义的递归定义 (Definition of Sheaf Semantics)}

设 $X$ 为 Locale，$U \in \mathcal{O}(X)$ 为开集。束语义通过递归定义 迫近关系（Forcing Relation） $U \models \varphi$。
这里 $\varphi$ 是作用在开集 $U$ 上的公式（其中的自由变量已被 $F(U)$ 中的截面 $s$ 所实例化）。

递归定义如下：

\begin{enumerate}
  \item 逻辑真 ($\top$)
  \[
  U \models \top \quad \iff \quad \text{总是成立 (true)}
  \]

  \item 逻辑假 ($\perp$)
  \[
  U \models \perp \quad \iff \quad U = \perp \text{（即 } U \text{ 是最小开集/空开集）}
  \]

  \item 等式 ($s = t$)
  对 $s, t \in F(U)$：
  \[
  U \models (s = t) \quad \iff \quad s = t \quad (\text{作为集合 } F(U) \text{ 中的元素完全相等})
  \]

  \item 合取 ($\wedge$)
  \[
  U \models \varphi \wedge \psi \quad \iff \quad (U \models \varphi) \text{ 且 } (U \models \psi)
  \]

  \item 析取 ($\vee$)
  \[
  U \models \varphi \vee \psi \quad \iff \quad \exists \text{ 覆盖 } U = \bigvee_{i \in I} U_i \text{ 使得对每个 } i \in I: (U_i \models \varphi) \text{ 或 } (U_i \models \psi)
  \]
  说明：束语义中的“或”不是全局的！不需要 $U$ 整体满足 $\varphi$ 或整体满足 $\psi$；只要能将 $U$ 局部打碎成开覆盖，在每个局部区域上满足其中之一即可。

  \item 蕴涵 ($\Rightarrow$)
  \[
  U \models \varphi \Rightarrow \psi \quad \iff \quad \forall V \le U, \quad \text{若 } V \models \varphi\vert{}_V, \text{ 则 } V \models \psi\vert{}_V
  \]
  说明：蕴涵是“向下遗传”的，必须对 $U$ 的所有子开集 $V$ 都成立。

  \item 存在量词 ($\exists$)
  设变量 $x$ 的类型为束 $F$：
  \[
  U \models \exists x : F. \, \varphi(x) \quad \iff \quad \exists \text{ 覆盖 } U = \bigvee_{i \in I} U_i, \, \text{且对每个 } i \in I, \exists s_i \in F(U_i), \, U_i \models \varphi(s_i)
  \]
  说明：存在性同样是局部的！不需要在整个 $U$ 上存在全局截面 $s \in F(U)$，只要能在局部开集 $U_i$ 上找到局部截面 $s_i$ 即可。

  \item 全称量词 ($\forall$)
  设变量 $x$ 的类型为束 $F$：
  \[
  U \models \forall x : F. \, \varphi(x) \quad \iff \quad \forall V \le U, \, \forall s_0 \in F(V), \quad V \models \varphi(s_0)
  \]
  说明：全称量词也是“向下遗传”的，必须检验 $U$ 的任意子开集 $V$ 上的每一个截面 $s_0 \in F(V)$。
\end{enumerate}

\section{束语义的基本定理与性质}

束语义满足以下关键的性质与定理，使得我们能够安全地在束的范畴内部进行逻辑推理。

\subsection{5.1 单调性定理 (Monotonicity)}

定理：设 $V \le U$。若 $U \models \varphi$，则 $V \models \varphi\vert{}_V$。

含义：如果开集 $U$ 迫近命题 $\varphi$，那么 $U$ 的任何子开集 $V$ 也迫近该命题。

\subsection{5.2 局部性定理 / 胶合性 (Locality)}

定理：设 $U = \bigvee_{i \in I} U_i$。若对所有 $i \in I$ 都有 $U_i \models \varphi\vert{}_{U_i}$，则 $U \models \varphi$。

含义：真值具有“束”的特征——如果命题在局部开覆盖的每一个块上都成立，则它在全局也成立。

\subsection{5.3 逻辑可靠性定理 (Soundness Theorem for Intuitionistic Logic)}

定理：若公式 $\varphi$ 在直觉主义一阶逻辑 (Intuitionistic First-Order Logic) 中可推导出 $\psi$（记作 $\varphi \vdash \psi$），且 $U \models \varphi$，则必然有 $U \models \psi$。

核心意义：束语义完全验证直觉主义逻辑！ 任何在构造性/直觉主义逻辑中成立的定理，在任何束范畴 $\text{Sh}(X)$ 的内部语义中都自动成立。但排中律 ($\varphi \vee \neg\varphi$) 在束语义中一般不成立。

\subsection{5.4 否定与双重否定 (Negation and Double Negation)}

在逻辑中，否定定义为 $\neg \varphi := (\varphi \Rightarrow \perp)$。

\begin{itemize}
  \item 否定的迫近：
  \[
  U \models \neg \varphi \quad \iff \quad \forall V \le U, \, (V \models \varphi \implies V = \perp)
  \]
  \item 双重否定的迫近 (Dense Open)：
  定理：$U \models \neg\neg \varphi$ 当且仅当存在一个在 $U$ 中稠密的开集 $V \le U$（即对任意非空 $W \le U$，$W \wedge V \ne \perp$），使得 $V \models \varphi$。
\end{itemize}

\section{具体应用与直观解释}

\subsection{6.1 分析学：连续函数束与介值定理 (Intermediate Value Theorem)}

设 $X$ 为拓扑空间，考虑实值连续函数束 $\mathcal{C}$，对开集 $U \subseteq X$，$\mathcal{C}(U) = \{ f: U \to \mathbb{R} \mid f \text{ 连续}\}$。

在 $\text{Sh}(X)$ 的束语义下，$\mathcal{C}$ 内部表现为“实数集” $\mathbb{R}$。

\begin{itemize}
  \item 经典介值定理失败的原因：
  设连续函数族 $f_x(t): \mathbb{R} \to \mathbb{R}$ 随参数 $x \in X$ 连续变化，且 $f_x(-1) < 0$ 且 $f_x(1) > 0$。
  内部命题 $\exists s: \mathcal{C}. \, f(s) = 0$ 成立，当且仅当存在开覆盖 $X = \bigcup_i U_i$，使得在每个 $U_i$ 上存在连续函数 $s: U_i \to \mathbb{R}$ 满足 $f_x(s(x)) = 0$。
  然而，对于某些连续函数族（例如在原点附近横穿但有重根或多分支），零点无法局部连续选择。因此，内部语言中经典的介值定理在 $\text{Sh}(X)$ 中不成立，这同时给出了“经典介值定理无法在构造性逻辑中证明”的语义解释。
  
  \item 构造性版本的成立：
  近似介值定理（$\forall \epsilon > 0, \exists x, \vert{}f(x)\vert{} < \epsilon$）和单调介值定理（严格单调函数零点唯一）在构造性逻辑中有证明，因此它们在任何 $\text{Sh}(X)$ 的束语义中自动成立。
\end{itemize}

\subsection{6.2 构造性代数：“任意化约环都是域” (Reduced Rings are Fields)}

在经典代数中，代数几何的极小素理想往往依赖选择公理（Zorn 补题）。
但在束语义中：
设 $A$ 为一个化约环（Reduced Ring，即 $x^n = 0 \Rightarrow x = 0$）。
可以在 $A$ 的谱 Locale $X = \text{Spec}(A)$ 上构造一个结构束 $A^\sim$。
在 $\text{Sh}(X)$ 的束语义下，$A^\sim$ 具有极佳的代数性质：

\begin{enumerate}
  \item $A^\sim$ 内部是一个“域 (Field)”：
  \[
  X \models \forall x : A^\sim. \, \left( \neg(\exists y : A^\sim. \, xy = 1) \Rightarrow x = 0 \right)
  \]
  
  \item 矩阵单射性质的无条件证明：
  对于化约环 $A$ 上的列数大于行数的矩阵 $M$，若 $M$ 在 $A$ 上单射，要证明 $1 = 0$。
  在传统代数中需要通过极小素理想局部化成域再用线性代数证明，这需要选择公理。
  但在束语义中，$M$ 保持作为 $A^\sim$ 上的单射矩阵。因为内部 $A^\sim$ 是域，由内部普通的线性代数定理直接得出矛盾，即内部 $X \models \perp$，从而解包外部直接得到 $A$ 中 $1 = 0$。
\end{enumerate}

\section{总结}

\begin{center}
\begin{tabular}{|p{4cm}|p{3.5cm}|p{7cm}|}
\hline
\textbf{概念/符号} & \textbf{对应数学实体} & \textbf{在束语义中的逻辑含义} \\
\hline
$X, \mathcal{O}(X)$ & Locale / 拓扑空间开集格 & 真值所处的结构（替代传统 $\{0,1\}$） \\
\hline
$F$ & Locale 上的束 & 内部语言中的某种“集合/类型” (Sort) \\
\hline
$s \in F(U)$ & 束 $F$ 在 $U$ 上的截面 & 局部参数化对象/项 \\
\hline
$U \models \varphi$ & 迫近关系 (Forcing) & 命题 $\varphi$ 在局部范围 $U$ 上验证为真 \\
\hline
$\vee, \exists$ & 析取与存在量词 & 具有局部性（只需局部开覆盖上成立） \\
\hline
$\Rightarrow, \forall$ & 蕴涵与全称量词 & 具有遗传性（对所有子开集/子截面成立） \\
\hline
\end{tabular}
\end{center}

束语义通过将“空间与束”翻译为“直觉主义逻辑内部模型”，搭建了一座连接拓扑学、范畴论与构造性代数/数理逻辑的桥梁。

\end{document}